% Extracted from AdaptiveAI(5).tex; manuscript snapshot, not preregistration.
\subsection{Search Strategy and Eligibility}

The Web of Science search targeted empirical
research on AI systems adapting to evolving tasks,
distributions, semantic categories, domains, or
operational environments. The 1,379 bibliographic
records retrieved constitute the fixed corpus of
this review.

Additional theoretical, methodological, and survey
references provide contextual coverage but are not
added to the primary-study corpus. Restricting the
search to a single bibliographic database limits
the coverage of the review and precludes claims
of exhaustive representation of adaptive AI research.

Study eligibility is defined through four inclusion
and four exclusion criteria, reflecting the scope,
adaptation characteristics, empirical evaluation,
and methodological interpretability of the
investigated systems.

\begin{description}
	\item[\textbf{IC1: Primary empirical research.}]
	Original empirical investigation of an AI or
	machine learning system.
	
	\item[\textbf{IC2: Identifiable adaptation.}]
	Modification of the system's learning state,
	parameters, representations, architecture,
	memory, or decision policy in response to
	changing conditions.
	
	\item[\textbf{IC3: Relevant evaluation.}]
	Measurable experimental outcomes associated
	with the adaptation mechanism.
	
	\item[\textbf{IC4: Methodological interpretability.}]
	Sufficient information to identify the adaptation
	problem, investigated mechanism, and evaluation
	conditions.
	
	\item[\textbf{EC1: Out-of-scope research.}]
	Research outside the operational definition
	established in Section~\ref{sec:introduction}.
	
	\item[\textbf{EC2: Non-empirical publication.}]
	Secondary research, editorials, perspectives,
	or conceptual contributions without relevant
	original empirical evaluation.
	
	\item[\textbf{EC3: Absence of adaptation.}]
	Static systems or systems without an
	identifiable adaptive mechanism.
	
	\item[\textbf{EC4: Absence of relevant evaluation.}]
	Studies without experimental evidence concerning
	the investigated adaptation mechanism.
\end{description}

Missing static comparators, retention measurements,
or resource assessments do not automatically render
an otherwise eligible study ineligible. Instead,
such omissions restrict the claims and research
questions that its evidence can address. This
distinction prevents the systematic removal of
incompletely reported studies whose methodological
limitations are themselves relevant to RQ2.
